

NCPC (Nordic Cargo Plane Control) are testing a new engine for their cargo planes. To this end they have bound a strong
and sturdy infinitely thin rope to the centre of their testing platform, and to the engine. We will place a coordinate system
onto this testing platform such that the rope is bound at the origin and lays along the positive $x$-axis to $(d, 0)$.
On this testing platform there are also a number of infinitely thin pillars that can stop the rope, but ignore the engine.
As the engine is started it starts rotating the rope counter-clockwise around the origin until it hits a pillar, at which point it 
is caught and starts rotating around that pillar counter-clockwise instead. The engine is then rotating at a smaller radius as some of
the rope is caught between the origin and this pillar. This keeps going until the rope is too short to reach any other pillars.

Running these tests, buying all this infinitely thin rope and setting up these infinitely thin pillars, is expensive. Besides, the
workers keep getting these nasty paper-like cuts from all these infinitely thin objects. It would be much more economical to just
simulate the behaviour. 

